\documentclass[12pt,a4paper]{article}

% --- Français ---
\usepackage[french]{babel}
\usepackage[utf8]{inputenc}
\usepackage[T1]{fontenc}
\usepackage{geometry}
\usepackage{xcolor}
\usepackage{listings}
\usepackage{hyperref}
\usepackage{fancyhdr}
\usepackage{graphicx}
\usepackage{setspace}
\usepackage{lmodern}

% --- Mise en page ---
\geometry{margin=2.2cm}
\setlength{\parindent}{0pt}
\setlength{\parskip}{5pt}
\onehalfspacing

% --- Couleurs ---
\definecolor{navy}{HTML}{1e293b}
\definecolor{indigo}{HTML}{6366f1}
\definecolor{emerald}{HTML}{059669}
\definecolor{codebg}{HTML}{f1f5f9}
\definecolor{codeborder}{HTML}{cbd5e1}
\definecolor{comment}{HTML}{6b7280}
\definecolor{keyword}{HTML}{4f46e5}
\definecolor{string}{HTML}{059669}

% --- Listings ---
\lstdefinestyle{phpphp}{
    language=PHP,
    backgroundcolor=\color{codebg},
    basicstyle=\small\ttfamily,
    keywordstyle=\color{keyword}\bfseries,
    stringstyle=\color{string},
    commentstyle=\color{comment}\itshape,
    numberstyle=\tiny\color{comment},
    numbers=left,
    numbersep=6pt,
    frame=single,
    rulecolor=\color{codeborder},
    breaklines=true,
    xleftmargin=10pt,
    framexleftmargin=6pt,
    tabsize=4,
    captionpos=b,
}
\lstdefinestyle{monsql}{
    language=SQL,
    backgroundcolor=\color{codebg},
    basicstyle=\small\ttfamily,
    keywordstyle=\color{keyword}\bfseries,
    stringstyle=\color{string},
    commentstyle=\color{comment}\itshape,
    numbers=left,
    numbersep=6pt,
    frame=single,
    rulecolor=\color{codeborder},
    breaklines=true,
    xleftmargin=10pt,
    framexleftmargin=6pt,
    tabsize=2,
    captionpos=b,
}

% --- En-tête ---
\pagestyle{fancy}
\fancyhf{}
\fancyhead[L]{\small\sffamily\textcolor{comment}{FasiChat Classroom}}
\fancyhead[R]{\small\sffamily\textcolor{comment}{Rapport de projet}}
\fancyfoot[C]{\thepage}
\renewcommand{\headrulewidth}{0.3pt}

% --- Liens ---
\hypersetup{
    colorlinks=true,
    linkcolor=indigo,
    urlcolor=indigo,
}

% --- Sections ---
\usepackage{titlesec}
\titleformat{\section}{\Large\sffamily\bfseries\color{navy}}{}{0em}{}[\vspace{-2pt}\rule{\textwidth}{0.4pt}]
\titleformat{\subsection}{\large\sffamily\bfseries\color{indigo!80!black}}{}{0em}{}

\begin{document}

% ===== PAGE DE GARDE =====
\begin{titlepage}
    \centering
    \vspace*{0.5cm}
    {\scshape\Large UNIVERSIT\'E PROTESTANTE AU CONGO \par}
    {\scshape\large FACULT\'E DES SCIENCES INFORMATIQUES \par}
    {\scshape\large KINSHASA/LINGWALA \par}
    \vspace{2cm}
    \vspace{1cm}
    {\scshape\Large Travail Pratique \par}
    \vspace{1cm}
    {\huge\bfseries FasiChatClassRoom \par}
    \vspace{2cm}

    \begin{flushleft}
    {\Large\bfseries Mati\`ere :} PROGRAMMATION WEB PHP \\
    \vspace{0.8cm}
    {\Large\bfseries Con\c cu par les \'etudiants :} \\
    {\large 1. KAGHOMA KITSIPA John} \\
    {\large 2. KIPAKA BAKARI Michel} \\
    {\large 3. MATHE MWERO Ben} \\
    \vspace{0.8cm}
    {\Large\bfseries Promotion :} L2 LMD \\
    \end{flushleft}

    \vfill
    {\large Ann\'ee Acad\'emique : 2025-2026 \par}
\end{titlepage}

% ===== TABLE DES MATIÈRES =====
\newpage
\tableofcontents
\newpage

% ===== 1. INTRODUCTION =====
\section{Introduction}

FasiChat est une application web qui permet aux \'etudiants, enseignants et
administrateurs d'une facult\'e de communiquer facilement. Le projet est n\'e
d'un besoin simple : remplacer les groupes WhatsApp et les annonces papier par
une plateforme unique, organis\'ee et s\'ecuris\'ee.

L'application permet :
\begin{itemize}
    \item d'\'echanger des messages priv\'es ou publics
    \item de publier des annonces officielles (Valve)
    \item d'afficher un mur p\'edagogique par cours
    \item d'envoyer des convocations aux enseignants
    \item de partager des fichiers (PDF, images, audio, vid\'eo)
\end{itemize}

\noindent Ce rapport d\'etaille l'architecture, les fonctionnalit\'es et les
choix techniques du projet.

% ===== 2. OBJECTIFS =====
\section{Objectifs}

L'objectif principal \'etait de cr\'eer un espace de communication unique pour
toute la facult\'e, avec des r\^oles bien d\'efinis :

\begin{itemize}
    \item \textbf{\'Etudiant :} peut discuter en groupe (promotion), voir le
    mur des cours, envoyer des messages priv\'es aux camarades, et consulter
    les annonces du Valve.
    \item \textbf{Enseignant / Assistant :} peut publier sur le mur
    p\'edagogique, discuter en public avec les \'etudiants, recevoir des
    convocations, et envoyer des messages priv\'es aux coll\`egues.
    \item \textbf{Doyen / Vice-Doyen :} peut convoquer des r\'eunions,
    \'echanger des messages confidentiels, et consulter tous les espaces.
    \item \textbf{Apparitaire :} g\`ere le tableau d'affichage (Valve) :
    publie, modifie et supprime les annonces officielles.
\end{itemize}

% ===== 3. TECHNOLOGIES =====
\section{Technologies utilis\'ees}

Le projet utilise uniquement des technologies natives, sans framework :

\begin{itemize}
    \item \textbf{PHP 8} c\^ot\'e serveur (POO, PDO, sessions)
    \item \textbf{MySQL} pour la base de donn\'ees
    \item \textbf{JavaScript} (Fetch API, DOM) pour le c\^ot\'e client
    \item \textbf{CSS3} (Flexbox, variables CSS, d\'egrad\'es)
    \item \textbf{Apache} (mod\_rewrite pour le routage)
    \item \textbf{HTML5} pour les templates
\end{itemize}

\noindent Le choix de ne pas utiliser de framework \'etait volontaire : cela
permet de mieux comprendre chaque brique du projet et de montrer notre
ma\^itrise des concepts fondamentaux.

% ===== 4. ARCHITECTURE =====
\section{Architecture du projet}

\subsection{Structure MVC}

Le projet suit une architecture MVC (Mod\`ele-Vue-Contr\^oleur) maison : les
\textbf{Mod\`eles} g\`erent les donn\'ees (Utilisateur, Message, Fichier\ldots),
les \textbf{Vues} affichent le HTML (login, dashboards, valve), et les
\textbf{Contr\^oleurs} traitent les requ\^etes.

\subsection{Front Controller}

Toutes les requ\^etes passent par \texttt{public/index.php}. Le fichier
\texttt{.htaccess} redirige les URLs vers ce fichier.

\begin{lstlisting}[style=phpphp, caption=Le routeur (public/index.php)]
$path = parse_url($_SERVER['REQUEST_URI'], PHP_URL_PATH);
$path = str_replace('/FasiChatClassroom/public', '', $path);
$path = trim($path, '/');

switch ($path) {
    case 'login':
        require __DIR__ . '/../views/login.php';
        break;
    case 'message-send':
        $msgModel = new Message($db);
        $success = $msgModel->envoyer(...);
        sendJsonResponse(['success' => $success]);
        break;
    default:
        require __DIR__ . '/../views/login.php';
}
\end{lstlisting}

\subsection{Hi\'erarchie des classes utilisateur}

Les utilisateurs sont repr\'esent\'es par une hi\'erarchie de classes PHP :
la classe de base \texttt{Utilisateur} est \textbf{abstraite} (on ne peut pas
cr\'eer un \guillemotleft~utilisateur~\guillemotright{} g\'en\'erique).
\texttt{CadreAdministratif} ajoute la m\'ethode \texttt{convoquer()} partag\'ee
par le Doyen et le Vice-Doyen.

\begin{lstlisting}[style=phpphp, caption=Classes utilisateur]
abstract class Utilisateur {
    protected $id, $nom, $prenom, $email, $role;
    public function __construct($data = []) { ... }
    public static function authentifier($db, $email, $password) { ... }
}
// Heritiers : Etudiant, Enseignant (et Assistant),
//             CadreAdministratif (abstrait),
//               - Doyen, ViceDoyen,
//             Apparitaire
\end{lstlisting}

% ===== 5. BASE DE DONNÉES =====
\section{Base de donn\'ees}

La base \texttt{fasichat} contient 7 tables reli\'ees par des cl\'es
\'etrang\`eres. Voici les plus importantes :

\begin{lstlisting}[style=monsql, caption=Tables principales]
-- Utilisateurs (7 roles)
CREATE TABLE utilisateurs (
    id INT AUTO_INCREMENT PRIMARY KEY,
    nom VARCHAR(100), prenom VARCHAR(100),
    email VARCHAR(100) UNIQUE,
    mot_de_passe VARCHAR(255), role VARCHAR(50)
);

-- Messages (type : 'prive', 'public', 'mur', 'convocation')
CREATE TABLE messages (
    id INT AUTO_INCREMENT PRIMARY KEY,
    expediteur_id INT,
    contenu TEXT,
    date_envoi DATETIME DEFAULT CURRENT_TIMESTAMP,
    type VARCHAR(20),
    cours_id INT NULL, promotion_id INT NULL,
    fichier_id INT NULL
);

-- Annonces (Valve)
CREATE TABLE annonces (
    id INT AUTO_INCREMENT PRIMARY KEY,
    titre VARCHAR(255), contenu TEXT,
    date_publication DATETIME DEFAULT CURRENT_TIMESTAMP,
    auteur_id INT
);
\end{lstlisting}

\noindent Les messages priv\'es sont li\'es \`a un destinataire via la table
\texttt{messages\_destinataires}. Les messages publics sont li\'es \`a une
promotion, les messages de mur sont li\'es \`a un cours.

% ===== 6. FONCTIONNALITÉS =====
\section{Fonctionnalit\'es d\'etaill\'ees}

\subsection{Messagerie en temps r\'eel}

Les messages sont envoy\'es et re\c{c}us sans rechargement de page. Toutes les
3 secondes, le navigateur demande les nouveaux messages au serveur (technique
de \textbf{polling AJAX}).

\begin{lstlisting}[style=phpphp, caption=Envoi d'un message (co^te' serveur)]
class Message {
    private $db;
    public function envoyer($expediteurId, $contenu, $type,
            $destinataireId = null, $coursId = null,
            $promotionId = null, $fichierId = null) {
        $this->db->beginTransaction();
        // Insertion du message
        $query = "INSERT INTO messages (...) VALUES (...)";
        $stmt->execute([...]);
        // Si prive, on lie le destinataire
        if ($type === 'prive') {
            $queryDest = "INSERT INTO messages_destinataires
                         (message_id, destinataire_id) VALUES (...)";
            $stmtDest->execute([...]);
        }
        $this->db->commit();
        return true;
    }
}
\end{lstlisting}

\subsection{Mur p\'edagogique}

Chaque cours a un mur o\`u l'enseignant publie des annonces, rappels ou
questions. Les \'etudiants voient ces publications dans leur sidebar. Les
messages de type \texttt{mur} sont filtr\'es par \texttt{cours\_id}.

\subsection{Valve (tableau d'affichage)}

Le Valve affiche toutes les annonces officielles, tri\'ees de la plus r\'ecente
\`a la plus ancienne. Seul l'Apparitaire peut publier, mais tout le monde peut
les consulter.

\begin{lstlisting}[style=phpphp, caption=Affichage des annonces]
$stmtAnnonces = $db->query("
    SELECT a.*, u.nom, u.prenom, u.role
    FROM annonces a
    JOIN utilisateurs u ON a.auteur_id = u.id
    ORDER BY a.date_publication DESC
");
$annonces = $stmtAnnonces->fetchAll();
\end{lstlisting}

\subsection{Convocations}

Le Doyen et le Vice-Doyen peuvent convoquer les enseignants. Le syst\`eme
cr\'ee un message de type \texttt{convocation} et une entr\'ee dans la table
\texttt{convocations} (objet, date, lieu).

\subsection{Upload de fichiers}

Le t\'el\'eversement se fait en deux \'etapes : (1) l'utilisateur
s\'electionne un fichier, le JavaScript l'envoie au serveur qui le stocke avec
un nom unique ; (2) le \texttt{fichier\_id} est joint au message lors de
l'envoi. Les images sont compress\'ees \`a 800px de largeur pour \'economiser
de l'espace.

\begin{lstlisting}[style=phpphp, caption=Upload se'curise' de fichier]
public function upload($fileArray) {
    $extension = strtolower(pathinfo($name, PATHINFO_EXTENSION));
    $safeName = uniqid('file_', true) . '.' . $extension;
    $destPath = __DIR__ . '/../../public/assets/uploads/' . $safeName;

    if (in_array($extension, ['jpg', 'jpeg', 'png'])) {
        $this->compressImage($fileArray['tmp_name'], $destPath, $ext);
    } else {
        move_uploaded_file($fileArray['tmp_name'], $destPath);
    }
    // Enregistrement en base...
    return $this->db->lastInsertId();
}
\end{lstlisting}

% ===== 7. SÉCURITÉ =====
\section{S\'ecurit\'e}

\begin{itemize}
    \item \textbf{Mots de passe :} hach\'es avec \texttt{password\_hash()}
    (bcrypt). V\'erification avec \texttt{password\_verify()}.
    \item \textbf{Injections SQL :} toutes les requ\^etes utilisent des
    requ\^etes pr\'epar\'ees PDO avec param\`etres nomm\'es.
    \item \textbf{XSS :} les donn\'ees affich\'ees sont \'echapp\'ees avec
    \texttt{htmlspecialchars()} c\^ot\'e PHP et \texttt{escapeHTML()} c\^ot\'e
    JavaScript.
    \item \textbf{Fichiers :} noms randomis\'es avec \texttt{uniqid()} pour
    \'eviter le path traversal.
    \item \textbf{Session :} chaque page v\'erifie la session et le r\^ole
    de l'utilisateur avant d'afficher le contenu.
\end{itemize}

% ===== 8. PROBLÈMES =====
\section{Probl\`emes rencontr\'es}

\subsection{Chemin d'upload incorrect}

\textbf{Probl\`eme :} les fichiers \'etaient stock\'es dans
\texttt{src/public/assets/uploads/} au lieu de
\texttt{public/assets/uploads/}. La constante \texttt{\_\_DIR\_\_} pointant
vers \texttt{src/Models/}, un seul \texttt{/../} remontait seulement \`a
\texttt{src/}.

\textbf{Solution :} utiliser \texttt{/../../} au lieu de \texttt{/../}.

\begin{lstlisting}[style=phpphp]
// Avant (faux)
$destPath = __DIR__ . '/../' . $uploadDir . $safeName;
// Apres (corrige)
$destPath = __DIR__ . '/../../' . $uploadDir . $safeName;
\end{lstlisting}

\subsection{Erreur de constructeur Doyen/Vice-Doyen}

\textbf{Probl\`eme :} \texttt{new Doyen(\$db, \$\_SESSION['user'])} passait
l'objet PDO comme donn\'ees utilisateur. Le constructeur
\texttt{Doyen(\$data)} interpr\'etait \texttt{\$db} comme \$data, ce qui
provoquait une erreur.

\textbf{Solution :} \texttt{new Doyen(\$\_SESSION['user'])} (un seul
param\`etre).

\subsection{Message avec fichier mais sans texte}

\textbf{Probl\`eme :} la validation exigeait un texte non vide. Si
l'utilisateur joignait un fichier sans \'ecrire de message, l'envoi \'echouait.

\textbf{Solution :} autoriser l'envoi quand un fichier est pr\'esent.

\begin{lstlisting}[style=phpphp]
// Avant
if (empty(trim($contenu))) { throw new Exception(...); }
// Apres
if (empty(trim($contenu)) && empty($fichierId)) { throw new Exception(...); }
\end{lstlisting}

\subsection{Mauvais type de messages}

\textbf{Probl\`eme :} les messages \'etudiant-enseignant \'etaient en
\texttt{type='prive'} au lieu de \texttt{type='public'}. Ils n'\'etaient
visibles que par deux personnes au lieu de toute la promotion.

\textbf{Solution :} modifier les appels JavaScript pour utiliser
\texttt{'public'} avec l'ID de la promotion.

\subsection{Absence de v\'erification d'auth}

\textbf{Probl\`eme :} les pages \texttt{dashboard\_apparitaire.php} et
\texttt{dashboard\_admin.php} n'avaient aucune v\'erification de session.

\textbf{Solution :} ajouter \texttt{session\_start()} et une redirection si
l'utilisateur n'est pas connect\'e.

% ===== 9. CONCLUSION =====
\section{Conclusion}

\subsection{Ce qui fonctionne}

Le projet r\'epond aux objectifs fix\'es :
\begin{itemize}
    \item messagerie temps r\'eel entre tous les acteurs
    \item distinction claire entre messages priv\'es, publics et de mur
    \item Valve (annonces) g\'er\'e par l'Apparitaire
    \item convocations pour les enseignants
    \item upload et compression de fichiers
    \item interface adapt\'ee \`a chaque r\^ole
    \item protection contre les injections SQL et XSS
\end{itemize}

\subsection{Pistes d'am\'elioration}

Quelques pistes pour aller plus loin :
\begin{itemize}
    \item remplacer le polling par des WebSockets (latence r\'eduite)
    \item ajouter des notifications (email ou push)
    \item interface mobile responsive d\'edi\'ee
    \item syst\`eme de \guillemotleft~lu~\guillemotright{} pour les messages
    \item administration des utilisateurs (CRUD)
\end{itemize}

\vspace{1cm}
\begin{center}
\rule{0.4\textwidth}{0.3pt}\par
\small\sffamily\textcolor{comment}{Fin du rapport}
\end{center}

\end{document}
